\documentclass[10pt,a4paper]{amsart}
\usepackage[margin=19mm]{geometry}
\usepackage{amsmath,amssymb,mathtools}
\usepackage[scr=rsfs,cal=euler]{mathalfa}
\usepackage{xfrac}
\usepackage[unicode,hidelinks,bookmarks=false]{hyperref}
\newcommand{\ie}{i.e.\ }
\newcommand{\cf}{cf.\ }
\newcommand{\resp}{respectively\ }
\newcommand{\Subsection}{\subsection}
\providecommand{\nobreakdash}{\nobreak\hbox{-}}
\setlength{\emergencystretch}{2em}
\multlinegap0pt
\usepackage[shortlabels]{enumitem}
\usepackage{amssymb}
\usepackage{mathtools}
\usepackage{xcolor}
\usepackage{algorithm}\usepackage{algpseudocode}
\newtheorem{lemma}{Lemma}
\DeclareMathOperator*{\E}{\mathbb{E}}
\newcommand{\Sph}{\mathbb{S}}
\newcommand{\cS}{\mathcal{S}}
\newcommand{\ip}[2]{\langle #1,#2\rangle}
\newcommand{\norm}[1]{\left\|#1\right\|}
\DeclareMathOperator{\range}{range}
\title{A Near-Optimal Lower Bound for $\ell_p$-Subspace Embeddings, $1\leq p<2$}
\author{Yi Li}
\date{Source: arXiv:2608.14201v1 (CC BY 4.0)}
\begin{document}
\maketitle

\noindent\emph{Source and licence.} A Near-Optimal Lower Bound for $\ell_p$-Subspace Embeddings, $1\leq p<2$. Version arXiv:2608.14201v1 (CC BY 4.0), distributed by arXiv under the Creative Commons Attribution 4.0 International licence, http://creativecommons.org/licenses/by/4.0/, as recorded in the arXiv licence metadata for this exact version and shown on the abstract page https://arxiv.org/abs/2608.14201v1. Full licence notice: \texttt{licenses/arxiv-2608.14201-CC-BY-4.0.txt}.\par\smallskip

\noindent\textit{Editorial scope.} Counted unit: the article's lemma lem:F\_properties (source0.tex lines 446--465). The article's complete original proof of it is delivered with it (source0.tex lines 466--502, one \texttt{proof} environment of 37 lines). No mathematics is written, repaired or re-derived: the statement and the proof are the article's own lines, in the article's own notation, and no retained line is substituted or re-flowed.  No further source-local context is needed: every cross-reference of every retained line resolves inside the two delivered spans.  Nothing was omitted from any retained span, so the package carries no omission record.  No retained line contains a citation, so the wrapper carries no bibliography and no bibliography entry.  No retained line contains a figure environment, an image-inclusion command or drawing code, so no image is delivered and the package declares no asset dependency.  Editorial formatting only, in addition: an AMS \texttt{amsart} preamble that restates the article's own declarations and macros used by the retained lines, namely the environments \texttt{lemma} on separate counters, together with the standard packages those lines need.  The retained environments print in this document as Lemma 1 in order of appearance, whereas the article prints the same items with the numbering of its own full document.  The complete native source of the article as distributed in the arXiv e-print for this exact version is bundled unchanged as \texttt{sources/207645/source0.tex}.
\begin{lemma}\label{lem:F_properties}
Suppose that $w\in \Sph^{2s-1}$ and $i\in I$ are arbitrary. The following properties hold.
\begin{enumerate}[label=(\roman*)]
	\item There exists an absolute constant $C_E > 0$ such that
	\[
		\left|	\frac{(2/p)\partial_\tau F(i,0,w) - b_0}{\eta\|R_i\|_2} 
		- w^\top \begin{pmatrix}
			        0 & W_i \\
			        W_i^\top & 0 \\
				 \end{pmatrix} w \right|
		\le C_E\delta,
	\]
	where $b_0 = (M\mathbf{1})_i$ is a constant independent of $i$.
	
	\item There is a constant $C_p > 0$ depending only on $p$ such that whenever $0\leq \tau \leq 1/4$,
	\[
		F(i,\tau,w)	\le C_p 2^k k^{p/2}.
	\]
\end{enumerate}
\end{lemma}

\begin{proof}
	Since $M$ and $\widetilde{M}$ agree on $\range(\widetilde{M})$, we have
	$\sum_{j\in U} M_{i,j} X_j = \sum_{j\in U}\widetilde{M}_{i,j}X_j = Y_i$. Therefore,
	\[
		\frac{\partial F}{\partial \tau}(i,0,w)
		= \frac{p}{2}\sum_{j\in U}|\langle i,j\rangle|^{p-2} w^\top C_j w 
		= \frac{p}{2}\left( \sum_{j\in U}|\langle i,j\rangle|^{p-2} + \eta w^\top \begin{pmatrix}
			0 & Y_i \\	Y_i^\top & 0 \\
		\end{pmatrix} w \right).
	\]
	Let $b_0 = \sum_{j\in U} |\ip{i}{j}|^{p-2}$, which is independent of $i$. We now have
	\[
		\frac{(2/p)\partial_\tau F(i,0,w) - b_0}{\eta\norm{R_i}_2} = w^\top \begin{pmatrix}
			0 & W_i \\	W_i^\top & 0 \\
		\end{pmatrix} w + \frac{1}{\norm{R_i}_2}
		w^\top \begin{pmatrix}
			0 & E_i \\	E_i^\top & 0 \\
		\end{pmatrix} w
	\]
	By the definition of $\cS$ and $\norm{R_i}_2\geq \sqrt{1-\delta^2}\norm{\widetilde{M}_i}_2$, we have
	\[
		\left| \frac{1}{\norm{R_i}_2}
			w^\top \begin{pmatrix}
				0 & E_i \\	E_i^\top & 0 \\
			\end{pmatrix} w \right| \leq C_E\delta.
	\]
	This proves (i).
	
	Next we prove (ii). Since $(3/4)I \preceq C_j\preceq (5/4)I$, we know that $w^\top C_j w \leq 5/4$. Hence
	\[
		F(i,\tau,w) \leq C_p\sum_{j\in U}\left(|\ip{i}{j}|^p+1\right).
	\]
	For uniformly random $j\in U$, $\ip{i}{j}$ has the distribution of a sum of $k$ independent Rademacher variables. Therefore $\E_j|\ip{i}{j}|^p\leq C_p k^{p/2}$ by the Khintchine inequality. It follows that
	\[
		F(i,\tau,w) \leq C_p2^k(k^{p/2}+1) \leq C_p2^k k^{p/2}. \qedhere
	\]
\end{proof}

\end{document}
